\bookStart{Fragments from Snorre’s Edda}
\setBookCode{EddicFragments}
%TODO: Further discussion on the fragments.

\section{Introduction}

A number of Eddic lines, stanzas and groups of stanzas are quoted in Snorre’s Edda.  The majority of them are taken from longer Eddic poems preserved in full in other manuscripts (primarily \Regius\ and \AM), but a few are found nowhere else.  These fragments will be edited in the present section.

The fragments have some things in common: they are generally pieces of spoken dialogue quoted in the context of longer narrative prose sections, and are, with one exception (Homedall’s galder, see below), not introduced by reference to their source but rather with phrases like \emph{þá kvað hann} ‘then he quoth’.

\sectionline

\section{I. A lost riddle-poem}\chapterStart%
\phantomsection\label{EddicFragments:F1}%

This half-stanza is quoted in \Gylfaginning~2.  It is the second Eddic stanza in \Gylfaginning, following \textlink{Havamal}[1] which is uttered by Yilfer himself when he enters the hall of the Eese in the same chapter.
The whole passage is clearly drawing on other Eddic mythic wisdom contests and is particularly reminiscent of \textlink{Vafthrudnismal}.  Whether the stanza itself is authentic or made up by Snorre himself is hard to tell.

\bpg\bpa%
\pva[15]Hann sá þrjú há-sę́ti ok hvert upp frá ǫðru, ok sátu þrír menn sinn í hverju.
\pva Þá spurði hann, hvert nafn hǫfðingja þeira vę́ri. Sá svarar, er hann leiddi inn, at
\pva sá, er í inu neðsta há-sę́ti sat, var konungr, ok heitir Hárr, en þar nę́st sá, er heitir Jafn·hárr, en sá ofast, er Þriði heitir.
\pva Þá spyrr Hárr komanda’nn, hvárt fleira er erendi hans, en heimill er matr ok drykkr honum sem ǫllum þar í Háva hǫll.
\pva Hann segir, at fyrst vill hann spyrja, ef nǫkkurr er fróðr maðr inni.
\pva Hárr segir, at hann komi eigi heill út, nema hann sé fróðari,\epa

\bpb\pvb[15]%
{\huge H}\textsc{e} [= Yilfer] \textsc{saw three high-seats}, and each one higher than the last; and three men sat there, each in his own seat.
\pvb Then he asked what the names of those chieftains were.  He who led him in answers that
\pvb the one who sat in the lowest high-seat was a king called High, and next to him he who is called Evenhigh, and uppermost he who is called Third.
\pvb Then High asks the guest whether he has any other errands---but he is as welcome to food and drink as every man is there in the High One’s hall.
\pvb He [= Yilfer] asks whether anyone within is a learned man.
\pvb High says that he will not come out whole unless he be the more learned one,\epb\epg

\bvg\bva%
\llap{„}ok statt-u \alst{f}ramm \hld\ meðan þú \alst{f}regn; &
\ind \alst{s}itja skal \alst{s}á es \alst{s}ęgir.“\eva

\bvb “and stand forth while thou askest; \\
\ind sit shall he who speaks!”\evb\evg

\sectionline

\section{II. Nearth and Shede}\chapterStart%
\phantomsection\label{EddicFragments:F2}%

The following passage is almost the whole of \Gylfaginning~23, excepting at the very end \emph{svá er sagt} ‘so it is said’, after which is quoted \textlink{Grimnismal}[11].
Notably, the two stanzas cited here are also found translated in \textcite{Saxo} 1.8.18--19, where they are said to have been spoken by the legendary hero Harding (corresponding to Nearth) and Rainhild (corresponding to Shede); in this way the old myth has been reinterpreted to fit Saxo’s heroic context.  For discussion see \textcite{Hopkins2021}.

\bpg\bpa%
Inn þriði ǫ́ss er sá, er kallaðr er Njǫrðr.  Hann býr á himni, þar sem heitir Nóa-tún.  Hann rę́ðr fyrir gǫngu vinds ok stillir sjá ok eld.  Á hann skal heita til sę́-fara ok til veiða.  Hann er svá auðigr ok fé-sę́ll, at hann má gefa þeim auð, landa eða lausa-fjár.  Á hann skal til þess heita.  \edtext{Eigi er Njǫrðr ása ę́ttar.  Hann var upp fǿddr í Vana-heimi, en Vanir gísluðu hann goðunum ok tóku í mót at gíslingu þann, er Hǿnir heitir.  Hann varð at sę́tt með goðum ok Vǫnum.}{\lemma{Eigi er Njǫrðr ása ę́ttar. \dots\ Hann varð at sę́tt með goðum ok Vǫnum. ‘Nearth is not of the lineage of the Eese. \dots\ He was used to reconcile the gods and the Wanes.’}\Bfootnote{For this episode see \textlink{Voluspa}[22] n., \textlink{Vafthrudnismal}[38]--39.}}  Njǫrðr á þá konu, er Skaði heitir, dóttir Þjatsa jǫtuns.  Skaði vill hafa bú-stað þann, er átt hafði faðir hennar, þat er á fjǫllum nǫkkurum, þar sem heitir Þrym-heimr, en Njǫrðr vill vera nę́r sę́.  Þau sę́ttust á þat, at þau skyldu vera níu nę́tr í Þrym-heimi, en þá aðrar níu at Nóa-túnum.  En er Njǫrðr kom aftr til Nóa-túna af fjalli’nu, þá kvað hann þetta:\epa

\bpb {\huge T}\textsc{he third Os is he who is called Nearth.}  He lives in the heaven in the place called Nowetons.  He rules the course of the wind, and stills sea and fire. On him shall one call for sea-faring and for hunting. He is so wealthy and blessed with money that he may give them a wealth of lands or loose property; on him shall one call for that sake. Nearth is not of the lineage of the Eese. He was brought up in Wanehome, but the Wanes gave him as a hostage to the gods, and in return got as hostage that one who is called Heener. He was used to reconcile the gods and the Wanes. Nearth has for wife the woman who is called Shede, the daughter of the ettin Thedse. Shede wishes to have the dwelling which her father had owned, which lies on some fells in the place called Thrimham---but Nearth wishes to live by the sea. They agreed with each other that they would live for nine nights in Thrimham, but the other nine at Nowetons. But when Nearth came back to the Nowetons from the fell, he quoth this:\epb\epg

\bvg\bva%
\llap{„}\alst{L}ęið eru⸗mk fjǫll, \hld\ vas’k⸗a \alst{l}ęngi ȧ, &
\ind \alst{n}ę́tr ęinar \alst{n}íu; &
\alst{u}lfa þytr \hld\ mér þȯtti \alst{i}llr vesa &
\ind hjá \alst{s}ǫngvi \alst{s}vana.“\eva

\bvb “Loathsome I find the fells, I was not long thereon--- \\
\ind only nine nights. \\
The howl of the wolves seemed ill to me \\
\ind against the song of swans.”\evb\evg

\bpg\bpa%
Þá kvað Skaði þetta:\epa

\bpb Then Shede quoth this:\epb\epg

\bvg\bva%
\llap{„}\alst{S}ofa né mát’k⸗a’k \hld\ \alst{s}ę́var bęðjum ȧ &
\ind \alst{f}ugls jarmi \alst{f}yrir; &
sá mik \alst{v}ękr \hld\ es af \alst{v}íði kømr &
\ind \alst{m}orgun hvęrjan \alst{m}ár.“\eva

\bvb “I could not sleep on the bed of the sea \\
\ind for the bleating of the bird. \\
That one awakens me who comes from the main \\
\ind every morning: the mew.”\evb\evg

\bpg\bpa%
Þá fór Skaði upp á fjall ok byggði í Þrym-heimi, ok ferr hón mjǫk á skíðum ok með boga ok skýtr dýr. Hón heitir ǫndur-goð eða ǫndur-dís.\epa

\bpb Then Shede went up to the fells and dwelled in Thrimham, and she often goes on skis with her bow and shoots beasts.  She is called ski-god or ski-dise.\epb\epg

\sectionline

\section{III. Homedall’s Galder (\emph{Hęim·dallar-galdr})}\chapterStart%
\phantomsection\label{EddicFragments:F3}%

This mysterious fragment is quoted in \Gylfaginning\ 27, the chapter describing Homedall, which is here reproduced in full.  The fragment consists of two c-lines and appears to be the end of a stanza in the fitting meter \Galdralag.

The same poem is mentioned again in \Skaldskaparmal\ 15: \emph{Heim·dallar hǫfuð heitir sverð. Svá er sagt, at hann var lostinn manns hǫfði í gegnum. Um þat er kveðit í Heim·dallar-galdri, ok er síðan kallat hǫfuð mjǫtuðr Heim·dallar} ‘A sword is called Homedall’s head. So it is said that he was pierced through with a man’s head; about that it is sung in Homedall’s galder, and thenceforth the head is called Homedall’s bane.’

\bpg\bpa%
\pva \edtrans{Heim·dallr}{Homedall}{\Bfootnote{“World-Brilliant”.  For the word \emph{dallr} ‘(possibly) brilliant, splendid’ see \textlink{Vafthrudnismal}[25] n.  Homedall appears to be an original fire-god, as seen by his epithets which resemble those of the Vedic \textoi{Agní}.}} heitir einn.
\pva Hann er kallaðr \edtrans{hvíti ǫ́ss}{the White Os}{\Bfootnote{Cf.~\textlink{Thrymskvida}[15]/1b: \emph{hvítastr ȧsa} ‘whitest of the Eese’.}};
\pva hann er mikill ok heilagr.
\pva Hann bǫ́ru at syni meyjar níu ok allar systr;
\pva hann heitir ok \edtrans{Hallin·skíði}{Haldenshid}{\Bfootnote{“The one with slanting sticks;” seemingly a name for a kind of (ritual?) fire.  It appears alongside \emph{Hęim·dali} as a poetic synonym for the ram in Þul \emph{Hrúts} 1 (\Skp\ 3).}} ok \edtrans{Gullin·tanni; tennr hans vǫ́ru af gulli}{Goldentooth; his teeth were of gold.}{\Bfootnote{The name is self-explanatory.  A direct cognate occurs in \Rigveda\ 5.2.3 as an adjective describing \textoi{Agní}, the god of fire: \textoi{ɦíraṇya-dantaṁ şúci-varṇam \dots\ a·paşyam} ‘I saw him with golden teeth and flaming color’, where OI \textoi{ɦíraṇya-danta} < PIE \emph{*ǵʰl̥‌h₃·enyó-h₃dontos} \char`~
\emph{*ǵʰl̥‌h₃t·iHno-h₃dontō} > PGmc \emph{*gulþīna-tanþô} > ON \emph{Gollin·tanni}.}}.
\pva Hestr hans heitir Gull·toppr.
\pva Hann býr þar er heitir Himin·bjǫrg við Bif·rǫst;
\pva hann er vǫrðr goða ok sitr þar við himins enda at gę́ta brúar’innar fyrir berg-risum.
\pva Hann þarf minna svefn en fugl.
\pva Hann sér jafnt nótt sem dag hundrað rasta frá sér;
\pva hann heyrir ok þat, er gras vex á jǫrðu eða ull á sauðum, ok allt þat er hę́ra lę́tr.
\pva Hann hefir lúðr þann er Gjallar-horn heitir, ok heyrir blástr hans í alla heima.
\pva Heim·dallar sverð er kallat hǫfuð manns.
\pva Hér er svá sagt: \edtext{[...]}{\Bfootnote{Here the text cites \textlink{Grimnismal}[13]; see there.}} \\
\pva Ok enn segir hann sjalfr í Heimdallar-galdri:\epa

\bpb\pvb {\huge H}\textsc{omedal is the name of one.}
\pvb He is called the White Os;
\pvb he is great and holy.
\pvb He was born as the son of nine maidens, sisters all.
\pvb He is also named Haldenshid and Goldentooth; his teeth were of gold.
\pvb His horse is named Goldtop.
\pvb He lives at the place called the Heavenbarrows near Bivrest.
\pvb He is the Watchman of the Gods and sits there at Heaven’s end to guard the bridge against barrow-risers.
\pvb He needs less sleep than a bird.
\pvb Both during night and day he sees as far as a hundred rests away from him;
\pvb he also hears when grass grows on the ground or wool on sheep, and everything which makes louder sound.
\pvb He has the lur called the Horn of Yell, and his blowing can be heard in all realms.
\pvb A man’s head is called ‘Homedall’s sword’.
\pvb Here it says so: [...] \\
\pvb And further he himself says in Homedall’s Galder:\epb\epg


\bvg\bva%
\llap{„}Níu em’k \edtrans{\alst{m}ęyja}{maidens}{\Afootnote{so \Upsaliensis; \emph{mǿðra} ‘mothers’ \RegiusProse\Trajectinus\Wormianus}\Bfootnote{I.e.~nine virgins.  The minority reading of \Upsaliensis\ is supported by Snorre’s prose in \Gylfaginning~27:4 which is clearly drawing from the present stanza.}} \alst{m}ǫgr, &
níu em’k \alst{s}ystra \edtrans{\alst{s}onr}{son}{\Afootnote{om. \Trajectinus}}.“\eva

\bvb “Of nine maidens am I the lad, \\
of nine sisters am I the son.”\evb\evg

\sectionline

\section{IV. Gna and the Wanes}\chapterStart%
\phantomsection\label{EddicFragments:F4}%

The following passage represents \Gylfaginning~35:21--26

\bpg\bpa\pva[21]%
Fjórtánda Gná, hana sendir Frigg í ymsa heima at ørindum sínum.
\pva Hón á þann hest, er renn lopt ok lǫg, er heitir Hóf·varpnir.
\pva Þat var eitt sinn, er hón reið, at vanir nǫkkvǫrir sá reið hennar í lopti’nu.
\pva Þá mę́lti einn:\epa

\bpb\pvb[21]%
{\huge T}\textsc{he fourteenth is Gna}; Frie sends her into various homes to run her errands.
\pvb She owns the horse which runs through air and sea and is called Hoofwarpner.
\pvb It was one time when she rode that some Wanes saw her riding in the air.
\pvb Then one spoke:\epb\epg


\bvg\bva%
\llap{„}Hvat þar \alst{f}lýgr, \hld\ hvat þar \alst{f}ęrr, &
\ind eða at \alst{l}opti \alst{l}íðr?“\eva

\bvb “What flies there, what fares there, \\
\ind or passes through the air?”\evb\evg


\bpg\bpa\pva[25]%
Hón svarar:\epa

\bpb\pvb[25]%
She answers:\epb\epg


\bvg\bva%
\llap{„}Né ek \alst{f}lýg, \hld\ þó ek \alst{f}ęr &
\ind ok at \alst{l}opti \alst{l}ið’k &
ȧ \alst{H}óf·varpni, \hld\ þęim’s \alst{H}am·skęrpir &
\ind \alst{g}at við \alst{G}arð·rofu.“\eva

\bvb “I fly not, though I fare \\
\ind and pass through the air, \\
on Hoofwarpner, whom Hamsherper \\
\ind begot with Yardrove.”\evb\evg


\bpg\bpa\pva[26]%
Af Gnár nafni er svá kallat, at þat gnę́far, er hǫ́tt ferr.\epa

\bpb\pvb[26]%
From Gna’s name it is so called that something which fares high up \emph{reaches high}.\epb\epg

\sectionline

\section{V. Balder’s death}\chapterStart%
\phantomsection\label{EddicFragments:F5}%

\Gylfaginning~49 contains the narrative of Balder’s death, beginning with his ominous dreams, and ending with the Eese failing to “weep him out of Hell” (for a summary and discussion of the myth and its attestations, see the introduction to \textlink{Voluspa}[31]--33). At the end of the chapter, a single \Ljodahattr\ speech-stanza is quoted.

The following excerpt represents \Gylfaginning~49:77--81, the very last part of that chapter.

\bpg\bpa\pva[77]%
Því nę́st sendu ę́sir um allan heim ørind-reka at biðja, at Baldr vę́ri grátinn ór Helju,
\pva en allir gerðu þat, menn’inir ok kykvendi’n ok jǫrð’in ok steinar’nir ok tré ok allr málmr,
\pva svá sem þú munt sét hafa, at þessir hlutir gráta, þá er þeir koma ór frosti ok í hita.
\pva Þá er sendi-menn fóru heim ok hǫfðu vel rekit sín ørindi, finna þeir í helli nǫkkvǫrum, hvar gýgr sat; hón nefndist Þǫkk.
\pva Þeir biðja hana gráta Baldr ór helju, hón segir:\epa

\bpb\pvb[77]%
{\huge T}\textsc{hereafter the Eese sent an errand-runner} through all the \inx[C]{Home} to ask that Balder be wept out of hell.
\pvb And all did that, the men and the beasts and the earth and the stones and trees and all metals,
\pvb even as thou must have seen that these things weep when they come out of cold and into heat.
\pvb When the messengers journeyed home and had run their errand well, they find in a certain cave where a \inx[C]{gow} was sitting; she called herself Thanks.
\pvb They ask her to weep Balder out of hell. She says:\epb\epg


\bvg\bva%
\llap{„}\alst{Þ}ǫkk mun gráta \hld\ \alst{þ}urrum tǫ́rum &
\ind \alst{B}aldrs \alst{b}ál-farar; &
\alst{k}yks né dauðs \hld\ naut’k⸗a \alst{K}arls sonar &
\ind \alst{h}afi \alst{H}ęl því’s \alst{h}ęfir.“\eva

\bvb “Thanks will weep dry tears \\
\ind for Balder’s pyre-journey \ken{death}. \\
Neither living nor dead I profitted from Churl’s son \ken*{= Balder}; \\
\ind let Hell have what she has!”\evb\evg


\bpg\bpa%
\pva[82]%
En þess geta menn, at þar hafi verit Loki Lauf·eyjar son, er flest hefir illt gørt með ǫ́sum.\epa

\bpb\pvb[82]%
But men guess that this must have been Lock, Leafy’s son, who has done the most evil among the Eese.\epb\epg

\sectionline

\section{VI. Thunder’s journey to Garfrith}\chapterStart%
\phantomsection\label{EddicFragments:F6}%

\Skaldskaparmal\ 26, here edited in part, tells of Thunder’s journey to the ettin Garfrith and his following killing of the ettin and his two daughters, Yelp and Grope.

Thunder’s fight with Garfrith and his daughters was clearly a very popular story in the late Wiking Age.  Its retelling takes up most of the Scaldic \Thorsdrapa, of which 19 full stanzas and 4 half-stanzas survive.  The killing of Yelp is also invoked in a pagan prayer to Thunder quoted in \Skaldskaparmal\ 11: \emph{stétt of Gjǫlp dauða} ‘Thou didst step over the dead Yelp’ (Vetrl Lv 1/1b in \Skp\ 3).  The story apparently remained popular after Christianization as seen by an following anecdote from the Kings’ Saws.  As king Harold ‘Hardrede’ Siwardson (1015--1066) was walking with his retinue, they witnessed an argument between a blacksmith and a tanner.  Harold asked one of his court poets to compose a stanza (viz. ÞjóðA Lv 5 in \Skp\ 2) about the incident while at the same time casting the blacksmith as Thunder and the tanner as Garfrith.  The king was impressed and asked the poet to make another stanza, this time about Siward and Fathomer.  The story of Garfrith is further alluded to in a curious way in \textcite[598--613]{Saxo} 8.14.1--20, where a party guided by a man named Thurkettle (Latin \emph{Thorkillus}) visit Garfrith’s (Latin \emph{Gerruthus}) cave and find his mutilated corpse along with those of three young women with their backs broken.  Thurkettle tells them that this was once done by the god Thunder (Latin \emph{Thor}).

Although the story of \Skaldskaparmal\ 26 is clearly based partly on \Thorsdrapa\ (from which it cites at length), it may also rely on a now-lost poem in \Ljodahattr\ since it quotes two stanzas in that meter.  Of these sts.~the first is found in all four main manuscripts while the second only in \Upsaliensis.  Both are spoken by Thunder and closely resemble each other stylistically, which is why they are likely to come from the same poem.

\bpg\bpa%
Þá fór Þórr til ár þeirar, er Vimur heitir, allra á mest.  Þá spennti hann sik megin-gjǫrðum ok studdi for·streymis Gríðar-vǫl, en Loki hélt undir megin-gjarðar.  Ok þá er Þórr kom á miðja ǫ́’na, þá óx svá mjǫk ǫ́’in, at uppi braut á ǫxl honum.  Þá kvað Þórr þetta:\epa

\bpb {\huge T}\textsc{hen journeyed Thunder} to that river which is called Wimbre, greatest of all rivers.  Then he strapped his might-girdle around him and leaned upon Grith’s stave against the stream, and Lock clung to the might-girdle.  And when Thunder came to the middle of the river, then it waxed so great that it broke over his shoulders. Then Thunder quoth this:\epb\epg


\bvg\bva%
\Ballnote{The contents of the st.~are also found in \Thorsdrapa\ 8/3--4.}%
\llap{„}\alst{V}ax⸗at-tu nú, \edtrans{\alst{V}imur}{Wimbre}{\Bfootnote{The name of the river is also found in a Thunder-kenning in \Husdrapa\ 6/3: \emph{Víð·gymnir Vimrar vaðs} ‘the Widegimner \name{ettin} of Wimbre’s ford’.}}, \hld\ alls mik þik \alst{v}aða tíðir &
\ind \alst{jǫ}tna garða \alst{ï}; &
\alst{v}ęitst, ef þú \alst{v}ęx \hld\ at þȧ \alst{v}ęx mér \edtrans{ǫ̇s-męgin}{Os-might}{\Bfootnote{TODO.}} &
\ind jafn-\alst{h}ǫ́tt upp sem \alst{h}iminn.“\eva

\bvb “Wax not now, Wimbre, as I wish to wade through thee \\
\ind into the yards of the ettins! \\
Thou knowest, if thou waxest, then my Os-might waxes \\
\ind up even as high as the heaven.”\evb\evg


\bpg\bpa%
Þá sér Þórr uppi í gljúfrum nǫkkurum, at Gjǫlp, dóttir Geir·røðar \edtrans{stóð þar tveim megin ár’innar, ok gerði hón ár-vǫxt’inn}{stood on both sides of the river, and she caused the river’s growth}{\Bfootnote{She stood with her legs spread and befouled the river with urine or menstrual blood.}}.  Þá tók Þórr upp ór ǫ́’nni stein mikinn ok kastaði at henni ok mę́lti svá: „At ósi skal ǫ́ stemma.“  Eigi missti hann, þar er hann kastaði til, ok í því bili bar hann at landi ok fekk tekit reyni-runn nǫkkurn ok steig svá ór ánni.  Því er þat orð-tak haft, at reynir er bjǫrg Þórs.\epa

\bpb Then Thunder sees that up in some certain gorges Yelp, daughter of Garfrith, stood on both sides of the river, and she caused the river’s growth.  Then Thunder took up from the river a great stone and threw it at her and spoke so: “At its source shall the river be dammed.” He did not miss his target, and in that moment he threw himself towards land and caught hold of a certain rowan shrub, and so he got out of the river.  From this comes the saying that the rowan is Thunder’s deliverance.\epb\epg


\bpg\bpa%
En er Þórr kom til Geir·røðar, þá var þeim fé-lǫgum vísat fyrst í geita-hús til her-bergis, ok var þar einn stóll til sę́tis, ok sat Þórr þar.  Þá varð hann þess varr, at stóll’inn fór undir honum upp at rę́fri.  Hann stakk Gríðar-veli upp í raft’ana ok lét sígast fast á stól’inn.  Varð þá brestr mikill, ok fylgði skrę́kr.  Þar hǫfðu verit undir stóli’num dǿtr Geir·røðar, Gjǫlp ok Greip, ok hafði hann brotit hrygg’inn í báðum.  \edtext{Þá kvað Þórr:}{\Afootnote{so \Upsaliensis; om. all others.}}\epa

\bpb But when Thunder came to Garfrith’s, then the companions were first shown into a goathouse for lodgings, and therein was but one chair for sitting and Thunder sat down there.  Then he became aware that the chair under him was moving upwards into the roof.  He thrust Grith’s stave up against the rafters and pushed himself firmly down into the chair.  Then there was a great crack, followed by a shriek.  There beneath the chair had been the daughters of Garfrith, Yelp and Grope, and he had broken the backs of them both.  Then Thunder quoth:\epb\epg

\bvg\bva%
\Aallnote{so \Upsaliensis; om. all others.}%
\llap{„}\alst{Ęi}nu \edtrans{\emph{sinni}}{time}{\Bfootnote{metr. and sens. emend.; om. \Upsaliensis.}} \hld\ nęytta’k \alst{a}lls męgins &
\ind \alst{jǫ}tna gǫrðum \alst{ï} &
þá’s \alst{G}jǫlp ok \alst{G}ręip, \hld\ dǿtr \alst{G}ęir-raðar, &
\ind vildu \alst{h}ęfja mik til \alst{h}imins.“\eva

\bvb “A single time I used all my might \\
\ind in the yards of the ettins, \\
when Yelp and Grope, the daughters of Garfrith, \\
\ind would lift me to the heaven.”\evb\evg

\sectionline

\section{VII. The tree Glazer}\chapterStart%
\phantomsection\label{EddicFragments:F7}%

The original context of this half-stanza in \Ljodahattr\ is obscure.  The present excerpt represents the whole ch. 42 of \Skaldskaparmal.

\bpg\bpa%
Hví er gull kallat barr eða lauf Glasis? Í Ǫ́s-garði fyrir durum Val-hallar stendr lundr, sá er Glasir er kallaðr, en lauf hans allt er gull-rautt, svá sem hér er kveðit, at\epa

\bpb {\huge W}\textsc{hy is gold} called the needle or leaf of Glazer? In Osyard, before the doors of Walhall stands a tree which is called Glazer, and its leafing is all golden red, as it is sung here, that\epb\epg

\bvg\bva%
\alst{G}lasir stendr \hld\ með \alst{g}ollnu laufi &
\ind fyrir \alst{S}igtýs \alst{s}ǫlum.\eva

\bvb Glazer stands with golden leaf \\
\ind before Sye-Tew’s \name{Weden’s} halls.\evb\evg

\bpg\bpa%
Sá er viðr fegrstr með goðum ok mǫnnum.\epa

\bpb It is the fairest tree among gods and men.\epb\epg

\sectionline

\section{VIII. On the making of Glapner}\chapterStart%
\phantomsection\label{EddicFragments:F8}%

The following fragmentary stanza about the making of Glapner---the fetter used to bind the Fenrerswolf---is found in the short work on kennings today called the \emph{Little Scalda} (\emph{Lítla skálda}), a text which probably served as a source for Snorre (for which see further \textcite[129--47]{Males2020}).

For the binding of the Wolf itself see \Gylfaginning~34:10--42 and the summary in \textlink{Lokasenna}[38]/3--4 n.

A form of this stanza is transparently paraphrased in \Gylfaginning~34:21, which tells what Glapner was made of: \emph{Hann var gǫrr af sex hlutum: af dyn kattar’ins ok af skeggi konu’nnar ok af rótum bjargs’ins ok af sinum bjarnar’ins ok af anda fisks’ins ok af fogls hráka.} ‘It was made from six things: from the cat’s din and from the woman’s beard and from the mountain’s root and from the bear’s sinews and from the fish’s breath and from the fowl’s spittle.’  The small variants---\emph{hráka} ‘spittle’ for \emph{mjǫlk} ‘milk’, and the reversal of lines 2 and 3---suggest that Snorre had access to a somewhat different variant of the present stanza.

\bvg\bva%
Ór \alst{k}attar dyn \hld\ ok ór \alst{k}onu skeggi, &
ór \alst{f}isks anda \hld\ ok ór \alst{f}ugla mjǫlk, &
ór \alst{b}ergs rótum \hld\ ok \alst{b}jarnar sinum, &
\ind ór því vas hann \alst{G}lęipnir \alst{g}ǫrr.\eva

\bvb “{\huge F}\textsc{rom cat’s din} and from woman’s beard, \\
from fish’s breath and from fowls’ milk, \\
from mountain’s roots and bear’s sinews--- \\
\ind from this was Glapner made.”\evb\evg

\sectionline
